\chapter{Dynamic programming II: masks, digits, ranges, and trees}\label{ch:dp2}
\index{bitmask DP}\index{digit DP}\index{range DP}\index{tree DP}

Chapter~11's recipe generalizes: states stop being positions and
become \emph{structured objects} --- subsets, automaton flags,
intervals, and tree-shape state pairs. Each generalization opens a
new problem class: bitmask DP for exponential-set problems, digit DP
for counting over enormous ranges, range DP for interval games, and
tree DP for hierarchical matching. After this chapter you will pick
the state shape, write the canonical template for each, and know
which optimizations wait at the horizon.

\section{Subsets as states}\label{sec:masks}\index{bitmask}
A subset of $\{1..n\}$ is an integer: mask DP stores one value per
subset, $O(2^n \cdot \text{poly})$ --- viable to $n \approx 20$.
Enumerating submasks uses \texttt{s = (s-1) \& mask}.
\begin{example}[CSES \emph{Elevator Rides}]\label{ex:elev}
Fewest rides given per-person weights and a lift capacity.
\par\smallskip
\textbf{Strategy.} State = the set already placed. dp[mask] tracks
\emph{two} things --- rides used and weight on the current ride ---
minimized lexicographically.
\par\smallskip
\begin{solution}
\lstinputlisting{code/elevator_rides.cpp}
Real output (CSES sample):
\begin{lstlisting}[numbers=none]
2
\end{lstlisting}
\end{solution}
\end{example}
\begin{figure}[ht]\centering
% fig55-masklattice.tex -- Figure 11.x: the subset lattice of {1,2,3,4}; a mask DP only moves upward (adds one element).
% fstyle.tex -- shared TikZ styles for the figures of this book (palette from lifeos.sty).
\tikzset{%
  cell/.style={draw=ink, minimum width=7mm, minimum height=6.5mm, inner sep=0pt, font=\footnotesize, fill=white},
  cellh/.style={cell, fill=accentsoft, draw=accent},
  cella/.style={cell, fill=amberbg, draw=amber},
  cellg/.style={cell, fill=goodbg, draw=good},
  cellx/.style={cell, fill=graybg, draw=slate, text=slate},
  vtx/.style={circle, draw=ink, minimum size=6.5mm, inner sep=0pt, font=\footnotesize, fill=white},
  vtxh/.style={vtx, draw=accent, fill=accentsoft},
  vtxa/.style={vtx, draw=amber, fill=amberbg},
  vtxg/.style={vtx, draw=good, fill=goodbg},
  vtxd/.style={vtx, draw=slate, fill=graybg, text=slate},
  lab/.style={font=\footnotesize, text=slate},
  labA/.style={font=\footnotesize, text=accent},
  labM/.style={font=\footnotesize, text=amber},
  idx/.style={font=\tiny, text=slate},
  arr/.style={-{stealth}, thick, ink},
  arrA/.style={-{stealth}, thick, accent},
  arrM/.style={-{stealth}, thick, amber},
  arrG/.style={-{stealth}, thick, good},
  edge/.style={thick, ink},
  edgeA/.style={thick, accent},
  edgeM/.style={thick, amber},
}%
% \cellrow{prefix}{x0}{y}{style}{v0,v1,...}: a row of cells named prefix0, prefix1, ... (pitch 0.75cm)
\providecommand{\cellrow}[5]{\foreach \v [count=\i from 0] in {#5} {\node[#4] (#1\i) at (#2+\i*0.75,#3) {\v};}}%
% \cellidx{prefix}{n}{start}: index labels under cells prefix0 .. prefix(n-1)
\providecommand{\cellidx}[3]{\foreach \i in {0,...,\numexpr#2-1\relax} {\node[idx, below=0pt of #1\i] {\the\numexpr\i+#3\relax};}}%
% \fitfig{<tikzpicture>}: shrink the picture to \linewidth only when it would overflow.
\ifdefined\fitbox\else\newsavebox{\fitbox}\fi%
\providecommand{\fitfig}[1]{\sbox{\fitbox}{#1}\ifdim\wd\fitbox>\linewidth\resizebox{\linewidth}{!}{\usebox{\fitbox}}\else\usebox{\fitbox}\fi}%

\fitfig{\begin{tikzpicture}[font=\footnotesize]
\node[vtxh, rounded corners=2pt, shape=rectangle, minimum width=9mm, minimum height=6mm, font=\scriptsize] (m0) at (0.00,0.00) {${\emptyset}$};
\node[vtx, rounded corners=2pt, shape=rectangle, minimum width=9mm, minimum height=6mm, font=\scriptsize] (m1) at (-2.17,1.25) {${1}$};
\node[vtx, rounded corners=2pt, shape=rectangle, minimum width=9mm, minimum height=6mm, font=\scriptsize] (m2) at (-0.72,1.25) {${2}$};
\node[vtx, rounded corners=2pt, shape=rectangle, minimum width=9mm, minimum height=6mm, font=\scriptsize] (m3) at (-3.62,2.50) {${12}$};
\node[vtx, rounded corners=2pt, shape=rectangle, minimum width=9mm, minimum height=6mm, font=\scriptsize] (m4) at (0.72,1.25) {${3}$};
\node[vtx, rounded corners=2pt, shape=rectangle, minimum width=9mm, minimum height=6mm, font=\scriptsize] (m5) at (-2.17,2.50) {${13}$};
\node[vtx, rounded corners=2pt, shape=rectangle, minimum width=9mm, minimum height=6mm, font=\scriptsize] (m6) at (-0.72,2.50) {${23}$};
\node[vtx, rounded corners=2pt, shape=rectangle, minimum width=9mm, minimum height=6mm, font=\scriptsize] (m7) at (-2.17,3.75) {${123}$};
\node[vtx, rounded corners=2pt, shape=rectangle, minimum width=9mm, minimum height=6mm, font=\scriptsize] (m8) at (2.17,1.25) {${4}$};
\node[vtx, rounded corners=2pt, shape=rectangle, minimum width=9mm, minimum height=6mm, font=\scriptsize] (m9) at (0.72,2.50) {${14}$};
\node[vtx, rounded corners=2pt, shape=rectangle, minimum width=9mm, minimum height=6mm, font=\scriptsize] (m10) at (2.17,2.50) {${24}$};
\node[vtx, rounded corners=2pt, shape=rectangle, minimum width=9mm, minimum height=6mm, font=\scriptsize] (m11) at (-0.72,3.75) {${124}$};
\node[vtx, rounded corners=2pt, shape=rectangle, minimum width=9mm, minimum height=6mm, font=\scriptsize] (m12) at (3.62,2.50) {${34}$};
\node[vtx, rounded corners=2pt, shape=rectangle, minimum width=9mm, minimum height=6mm, font=\scriptsize] (m13) at (0.72,3.75) {${134}$};
\node[vtx, rounded corners=2pt, shape=rectangle, minimum width=9mm, minimum height=6mm, font=\scriptsize] (m14) at (2.17,3.75) {${234}$};
\node[vtxh, rounded corners=2pt, shape=rectangle, minimum width=9mm, minimum height=6mm, font=\scriptsize] (m15) at (0.00,5.00) {${1234}$};
\draw[slate] (m0.north) -- (m1.south);
\draw[slate] (m0.north) -- (m2.south);
\draw[slate] (m0.north) -- (m4.south);
\draw[slate] (m0.north) -- (m8.south);
\draw[slate] (m1.north) -- (m3.south);
\draw[slate] (m1.north) -- (m5.south);
\draw[slate] (m1.north) -- (m9.south);
\draw[slate] (m2.north) -- (m3.south);
\draw[slate] (m2.north) -- (m6.south);
\draw[slate] (m2.north) -- (m10.south);
\draw[slate] (m3.north) -- (m7.south);
\draw[slate] (m3.north) -- (m11.south);
\draw[slate] (m4.north) -- (m5.south);
\draw[slate] (m4.north) -- (m6.south);
\draw[slate] (m4.north) -- (m12.south);
\draw[slate] (m5.north) -- (m7.south);
\draw[slate] (m5.north) -- (m13.south);
\draw[slate] (m6.north) -- (m7.south);
\draw[slate] (m6.north) -- (m14.south);
\draw[slate] (m7.north) -- (m15.south);
\draw[slate] (m8.north) -- (m9.south);
\draw[slate] (m8.north) -- (m10.south);
\draw[slate] (m8.north) -- (m12.south);
\draw[slate] (m9.north) -- (m11.south);
\draw[slate] (m9.north) -- (m13.south);
\draw[slate] (m10.north) -- (m11.south);
\draw[slate] (m10.north) -- (m14.south);
\draw[slate] (m11.north) -- (m15.south);
\draw[slate] (m12.north) -- (m13.south);
\draw[slate] (m12.north) -- (m14.south);
\draw[slate] (m13.north) -- (m15.south);
\draw[slate] (m14.north) -- (m15.south);
\foreach \k/\t in {0/{size 0},1/{size 1},2/{size 2},3/{size 3},4/{size 4}} \node[idx, anchor=west] at (5.6,\k*1.25) {\t};
\node[labA, anchor=west, align=left] at (6.7,2.5) {a transition adds one element:\\ the DP climbs from $\emptyset$ to the full set};
\end{tikzpicture}}

\caption{The subset lattice. Notice a mask DP always walks
upward --- each transition adds exactly one element.}
\label{fig:mask}
\end{figure}

\section{Digit DP}\label{sec:digit}\index{digit DP}
Counting integers $\le N$ with a digit property: walk the digits
most-significant-first carrying a \emph{tight} flag (still matching
$N$'s prefix?) and a \emph{leading-zero} flag (the number hasn't
started). Tight bounds the next digit; loose reads from memo.
\begin{example}[CSES \emph{Counting Numbers}]\label{ex:num}
Count integers in $[a,b]$ with no two adjacent digits equal.
\par\smallskip
\textbf{Strategy.} The automaton is (pos, tight, leading,
previous-digit); count(b) $-$ count(a$-$1).
\par\smallskip
\begin{solution}
\lstinputlisting{code/counting_numbers.cpp}
Real output (CSES sample, $a{=}123$, $b{=}321$):
\begin{lstlisting}[numbers=none]
171
\end{lstlisting}
\end{solution}
\end{example}
\begin{figure}[ht]\centering
% fig56-digitauto.tex -- Figure 11.x: tight/loose automaton of a digit DP (bound = 321, adjacent digits must differ).
% fstyle.tex -- shared TikZ styles for the figures of this book (palette from lifeos.sty).
\tikzset{%
  cell/.style={draw=ink, minimum width=7mm, minimum height=6.5mm, inner sep=0pt, font=\footnotesize, fill=white},
  cellh/.style={cell, fill=accentsoft, draw=accent},
  cella/.style={cell, fill=amberbg, draw=amber},
  cellg/.style={cell, fill=goodbg, draw=good},
  cellx/.style={cell, fill=graybg, draw=slate, text=slate},
  vtx/.style={circle, draw=ink, minimum size=6.5mm, inner sep=0pt, font=\footnotesize, fill=white},
  vtxh/.style={vtx, draw=accent, fill=accentsoft},
  vtxa/.style={vtx, draw=amber, fill=amberbg},
  vtxg/.style={vtx, draw=good, fill=goodbg},
  vtxd/.style={vtx, draw=slate, fill=graybg, text=slate},
  lab/.style={font=\footnotesize, text=slate},
  labA/.style={font=\footnotesize, text=accent},
  labM/.style={font=\footnotesize, text=amber},
  idx/.style={font=\tiny, text=slate},
  arr/.style={-{stealth}, thick, ink},
  arrA/.style={-{stealth}, thick, accent},
  arrM/.style={-{stealth}, thick, amber},
  arrG/.style={-{stealth}, thick, good},
  edge/.style={thick, ink},
  edgeA/.style={thick, accent},
  edgeM/.style={thick, amber},
}%
% \cellrow{prefix}{x0}{y}{style}{v0,v1,...}: a row of cells named prefix0, prefix1, ... (pitch 0.75cm)
\providecommand{\cellrow}[5]{\foreach \v [count=\i from 0] in {#5} {\node[#4] (#1\i) at (#2+\i*0.75,#3) {\v};}}%
% \cellidx{prefix}{n}{start}: index labels under cells prefix0 .. prefix(n-1)
\providecommand{\cellidx}[3]{\foreach \i in {0,...,\numexpr#2-1\relax} {\node[idx, below=0pt of #1\i] {\the\numexpr\i+#3\relax};}}%
% \fitfig{<tikzpicture>}: shrink the picture to \linewidth only when it would overflow.
\ifdefined\fitbox\else\newsavebox{\fitbox}\fi%
\providecommand{\fitfig}[1]{\sbox{\fitbox}{#1}\ifdim\wd\fitbox>\linewidth\resizebox{\linewidth}{!}{\usebox{\fitbox}}\else\usebox{\fitbox}\fi}%

\fitfig{\begin{tikzpicture}[font=\footnotesize]
\node[cellh, minimum width=2.2cm, minimum height=1.1cm, rounded corners=3pt, font=\small] (t) at (0,0) {tight};
\node[cell, minimum width=2.2cm, minimum height=1.1cm, rounded corners=3pt, font=\small] (l) at (6,0) {loose};
\draw[arrM, bend left=22] (t) to node[above, labM, align=center] {digit $<$ bound digit} (l);
\draw[arrA] (t) to[out=150,in=210,looseness=6] node[left, labA, align=right] {digit $=$ bound digit} (t);
\draw[arr] (l) to[out=30,in=-30,looseness=6] node[right, lab, align=left] {any digit\\ $\ne$ previous} (l);
\node[lab, align=center] at (3,-1.6) {tight never comes back: one undercut frees the whole suffix, and loose states are memoised};
\node[lab, align=center] at (3,-2.3) {bound $321$: tight loop reads $3,2,1$; picking $1<2$ in the second place jumps to loose};
\end{tikzpicture}}

\caption{The tight/loose automaton. Notice tight never comes back ---
one undercut frees the whole suffix.}
\label{fig:digitauto}
\end{figure}

\section{Range DP}\label{sec:range}\index{range DP}
Cost over every subarray: dp[l][r] built \emph{by length}, since a
range needs its strictly-smaller subranges first. Splits and
endpoint choices both fit the shape.
\begin{example}[CSES \emph{Removal Game}]\label{ex:removal}
Two players pick ends; maximize the first player's sum.
\par\smallskip
\textbf{Strategy.} dp[l][r] = best score difference from this range:
$\max(a_l - dp[l{+}1][r], a_r - dp[l][r{-}1])$ --- opponent-optimal
play subtracts.
\par\smallskip
\begin{solution}
\lstinputlisting{code/removal_game.cpp}
Real output (CSES sample):
\begin{lstlisting}[numbers=none]
8
\end{lstlisting}
\end{solution}
\end{example}
\begin{figure}[ht]\centering
% fig57-rangetri.tex -- Figure 11.x: dp[l][r] = best score difference on a[l..r] for a = {4,5,1,3} (Removal Game sample); answer (13+3)/2 = 8.
% fstyle.tex -- shared TikZ styles for the figures of this book (palette from lifeos.sty).
\tikzset{%
  cell/.style={draw=ink, minimum width=7mm, minimum height=6.5mm, inner sep=0pt, font=\footnotesize, fill=white},
  cellh/.style={cell, fill=accentsoft, draw=accent},
  cella/.style={cell, fill=amberbg, draw=amber},
  cellg/.style={cell, fill=goodbg, draw=good},
  cellx/.style={cell, fill=graybg, draw=slate, text=slate},
  vtx/.style={circle, draw=ink, minimum size=6.5mm, inner sep=0pt, font=\footnotesize, fill=white},
  vtxh/.style={vtx, draw=accent, fill=accentsoft},
  vtxa/.style={vtx, draw=amber, fill=amberbg},
  vtxg/.style={vtx, draw=good, fill=goodbg},
  vtxd/.style={vtx, draw=slate, fill=graybg, text=slate},
  lab/.style={font=\footnotesize, text=slate},
  labA/.style={font=\footnotesize, text=accent},
  labM/.style={font=\footnotesize, text=amber},
  idx/.style={font=\tiny, text=slate},
  arr/.style={-{stealth}, thick, ink},
  arrA/.style={-{stealth}, thick, accent},
  arrM/.style={-{stealth}, thick, amber},
  arrG/.style={-{stealth}, thick, good},
  edge/.style={thick, ink},
  edgeA/.style={thick, accent},
  edgeM/.style={thick, amber},
}%
% \cellrow{prefix}{x0}{y}{style}{v0,v1,...}: a row of cells named prefix0, prefix1, ... (pitch 0.75cm)
\providecommand{\cellrow}[5]{\foreach \v [count=\i from 0] in {#5} {\node[#4] (#1\i) at (#2+\i*0.75,#3) {\v};}}%
% \cellidx{prefix}{n}{start}: index labels under cells prefix0 .. prefix(n-1)
\providecommand{\cellidx}[3]{\foreach \i in {0,...,\numexpr#2-1\relax} {\node[idx, below=0pt of #1\i] {\the\numexpr\i+#3\relax};}}%
% \fitfig{<tikzpicture>}: shrink the picture to \linewidth only when it would overflow.
\ifdefined\fitbox\else\newsavebox{\fitbox}\fi%
\providecommand{\fitfig}[1]{\sbox{\fitbox}{#1}\ifdim\wd\fitbox>\linewidth\resizebox{\linewidth}{!}{\usebox{\fitbox}}\else\usebox{\fitbox}\fi}%

\begin{tikzpicture}[font=\footnotesize]
\tikzset{mc/.style={draw=slate, minimum width=9mm, minimum height=7.5mm, inner sep=0pt}}
\node[idx] at (0.0,0.65) {$r=1$};
\node[idx] at (1.0,0.65) {$r=2$};
\node[idx] at (2.0,0.65) {$r=3$};
\node[idx] at (3.0,0.65) {$r=4$};
\node[idx, anchor=east] at (-0.55,0.00) {$l=1$};
\node[idx, anchor=east] at (-0.55,-0.85) {$l=2$};
\node[idx, anchor=east] at (-0.55,-1.70) {$l=3$};
\node[idx, anchor=east] at (-0.55,-2.55) {$l=4$};
\node[mc, fill=accent!25, draw=accent] (c00) at (0.0,0.00) {4};
\node[mc, fill=accentsoft] (c01) at (1.0,0.00) {1};
\node[mc, fill=accentsoft] (c02) at (2.0,0.00) {0};
\node[mc, fill=amber!25, draw=amber] (c03) at (3.0,0.00) {3};
\node[mc, fill=accent!25, draw=accent] (c11) at (1.0,-0.85) {5};
\node[mc, fill=accentsoft] (c12) at (2.0,-0.85) {4};
\node[mc, fill=accentsoft] (c13) at (3.0,-0.85) {3};
\node[mc, fill=accent!25, draw=accent] (c22) at (2.0,-1.70) {1};
\node[mc, fill=accentsoft] (c23) at (3.0,-1.70) {2};
\node[mc, fill=accent!25, draw=accent] (c33) at (3.0,-2.55) {3};
\draw[arrM, shorten >=1pt] (c11.east) -- (c12.west);\draw[arrM, shorten >=1pt] (c22.north) -- (c12.south);
\node[labM, anchor=west, align=left] at (4.0,-0.4) {$dp[l][r]=\max\big(a_l-dp[l{+}1][r],\ a_r-dp[l][r{-}1]\big)$};
\node[lab, anchor=west, align=left, text width=7.5cm] at (4.0,-1.6) {diagonal first (length 1, blue), then longer ranges: every cell reads two shorter ranges already computed (arrows show $dp[2][3]$)};
\node[labM, anchor=west] at (4.0,-2.8) {$dp[1][4]=3$, so the first player gets $(13+3)/2=8$};
\end{tikzpicture}

\caption{By-length filling on the triangle. Notice the diagonal is
free --- every other cell reads smaller ranges already computed.}
\label{fig:rangetri}
\end{figure}
\begin{example}[CSES \emph{Edit Distance}]\label{ex:edit}
Minimum edits between two strings.
\par\smallskip
\textbf{Strategy.} dp[i][j] aligns prefixes: diagonal for
match/substitute, left and up for insert/delete. LCS is the same
grid with only diagonal-or-skip moves.
\par\smallskip
\begin{solution}
\lstinputlisting{code/edit_distance.cpp}
Real output (CSES sample):
\begin{lstlisting}[numbers=none]
2
\end{lstlisting}
\end{solution}
\end{example}
\begin{figure}[ht]\centering
\begin{tikzpicture}[font=\tiny,x=0.45cm,y=0.45cm]
\foreach \x in {0,...,4} \foreach \y in {0,...,5}
  \draw[slate!50] (\x,\y) rectangle +(1,1);
\draw[-{stealth},amber,thick] (0.5,4.5)--(1.5,3.5)--(2.5,2.5)--(3.5,1.5)--(3.5,0.5);
\node[font=\tiny] at (2,-0.8) {diagonal = match/substitute};
\end{tikzpicture}

\caption{The alignment grid with its optimal path. Notice every
path is an alignment --- the DP is just shortest path on this DAG.}
\label{fig:editpath}
\end{figure}

\section{DP on trees}\label{sec:treedp}\index{tree DP}
Rooted trees carry state pairs: \emph{in}[v] = best when $v$ is
matched downward, \emph{out}[v] = best when $v$ is free to match
upward. Chapter~10's reroot machinery then upgrades per-root answers
to all-roots.
\begin{example}[CSES \emph{Tree Matching}]\label{ex:tmatch}
Maximum edges with no shared endpoint.
\par\smallskip
\textbf{Strategy.} in/out states: \texttt{out[v]} sums children's
max; \texttt{in[v]} picks the child offering the best edge gain.
\par\smallskip
\begin{solution}
\lstinputlisting{code/tree_matching.cpp}
Real output (5-vertex demo):
\begin{lstlisting}[numbers=none]
2
\end{lstlisting}
\end{solution}
\end{example}
\begin{figure}[ht]\centering
\begin{tikzpicture}[font=\footnotesize]
\node[circle,draw=slate,font=\tiny] (v) at (0,0.9) {v};
\node[circle,draw=slate,font=\tiny] (a) at (-0.7,0) {a};
\node[circle,draw=slate,font=\tiny] (b) at (0.7,0) {b};
\draw[slate] (v)--(a) (v)--(b);
\node[font=\tiny,anchor=west] at (1.4,0.9) {in[v]: v matched below};
\node[font=\tiny,anchor=west] at (1.4,0.3) {out[v]: v free upward};
\node[font=\tiny] at (0.7,-0.8) {two states per vertex};
\end{tikzpicture}

\caption{The in/out pair. Notice a vertex's ``free upward'' is what
the parent's in-state needs --- the pair is the interface.}
\label{fig:treestate}
\end{figure}

\section{The optimisation horizon}\label{sec:horizon}\index{bitset}
One optimization worth teaching now: \emph{bitset knapsack} ---
reachable sums as a bitmask, \texttt{dp |= dp << w} --- $O(nC/64)$.
\begin{example}[CSES \emph{Money Sums}]\label{ex:money}
All reachable coin sums.
\par\smallskip
\textbf{Strategy.} A sum is a set member: shift the bitset by each
coin and OR. One machine word covers 64 reachable positions.
\par\smallskip
\begin{solution}
\lstinputlisting{code/money_sums.cpp}
Real output (CSES sample):
\begin{lstlisting}[numbers=none]
6
2 3 5 7 8 10
\end{lstlisting}
\end{solution}
\end{example}
\begin{figure}[ht]\centering
% fig60-bitset.tex -- Figure 11.x: Money Sums as bitset shift-or with coins {2,3,5}; reachable sums 2 3 5 7 8 10.
% fstyle.tex -- shared TikZ styles for the figures of this book (palette from lifeos.sty).
\tikzset{%
  cell/.style={draw=ink, minimum width=7mm, minimum height=6.5mm, inner sep=0pt, font=\footnotesize, fill=white},
  cellh/.style={cell, fill=accentsoft, draw=accent},
  cella/.style={cell, fill=amberbg, draw=amber},
  cellg/.style={cell, fill=goodbg, draw=good},
  cellx/.style={cell, fill=graybg, draw=slate, text=slate},
  vtx/.style={circle, draw=ink, minimum size=6.5mm, inner sep=0pt, font=\footnotesize, fill=white},
  vtxh/.style={vtx, draw=accent, fill=accentsoft},
  vtxa/.style={vtx, draw=amber, fill=amberbg},
  vtxg/.style={vtx, draw=good, fill=goodbg},
  vtxd/.style={vtx, draw=slate, fill=graybg, text=slate},
  lab/.style={font=\footnotesize, text=slate},
  labA/.style={font=\footnotesize, text=accent},
  labM/.style={font=\footnotesize, text=amber},
  idx/.style={font=\tiny, text=slate},
  arr/.style={-{stealth}, thick, ink},
  arrA/.style={-{stealth}, thick, accent},
  arrM/.style={-{stealth}, thick, amber},
  arrG/.style={-{stealth}, thick, good},
  edge/.style={thick, ink},
  edgeA/.style={thick, accent},
  edgeM/.style={thick, amber},
}%
% \cellrow{prefix}{x0}{y}{style}{v0,v1,...}: a row of cells named prefix0, prefix1, ... (pitch 0.75cm)
\providecommand{\cellrow}[5]{\foreach \v [count=\i from 0] in {#5} {\node[#4] (#1\i) at (#2+\i*0.75,#3) {\v};}}%
% \cellidx{prefix}{n}{start}: index labels under cells prefix0 .. prefix(n-1)
\providecommand{\cellidx}[3]{\foreach \i in {0,...,\numexpr#2-1\relax} {\node[idx, below=0pt of #1\i] {\the\numexpr\i+#3\relax};}}%
% \fitfig{<tikzpicture>}: shrink the picture to \linewidth only when it would overflow.
\ifdefined\fitbox\else\newsavebox{\fitbox}\fi%
\providecommand{\fitfig}[1]{\sbox{\fitbox}{#1}\ifdim\wd\fitbox>\linewidth\resizebox{\linewidth}{!}{\usebox{\fitbox}}\else\usebox{\fitbox}\fi}%

\begin{tikzpicture}[font=\footnotesize]
\tikzset{mc/.style={draw=slate, minimum width=6.4mm, minimum height=6mm, inner sep=0pt, font=\scriptsize}}
\node[lab, anchor=east] at (-0.5,0.6) {sum};
\node[idx] at (0.00,0.6) {0};
\node[idx] at (0.68,0.6) {1};
\node[idx] at (1.36,0.6) {2};
\node[idx] at (2.04,0.6) {3};
\node[idx] at (2.72,0.6) {4};
\node[idx] at (3.40,0.6) {5};
\node[idx] at (4.08,0.6) {6};
\node[idx] at (4.76,0.6) {7};
\node[idx] at (5.44,0.6) {8};
\node[idx] at (6.12,0.6) {9};
\node[idx] at (6.80,0.6) {10};
\node[lab, anchor=east] at (-0.5,0.00) {start};
\node[mc, fill=accentsoft, draw=accent] at (0.00,0.00) {1};
\node[mc, text=slate] at (0.68,0.00) {0};
\node[mc, text=slate] at (1.36,0.00) {0};
\node[mc, text=slate] at (2.04,0.00) {0};
\node[mc, text=slate] at (2.72,0.00) {0};
\node[mc, text=slate] at (3.40,0.00) {0};
\node[mc, text=slate] at (4.08,0.00) {0};
\node[mc, text=slate] at (4.76,0.00) {0};
\node[mc, text=slate] at (5.44,0.00) {0};
\node[mc, text=slate] at (6.12,0.00) {0};
\node[mc, text=slate] at (6.80,0.00) {0};
\node[lab, anchor=east] at (-0.5,-0.62) {after coin 2};
\node[mc, fill=accentsoft, draw=accent] at (0.00,-0.62) {1};
\node[mc, text=slate] at (0.68,-0.62) {0};
\node[mc, fill=amberbg, draw=amber, text=amber] at (1.36,-0.62) {1};
\node[mc, text=slate] at (2.04,-0.62) {0};
\node[mc, text=slate] at (2.72,-0.62) {0};
\node[mc, text=slate] at (3.40,-0.62) {0};
\node[mc, text=slate] at (4.08,-0.62) {0};
\node[mc, text=slate] at (4.76,-0.62) {0};
\node[mc, text=slate] at (5.44,-0.62) {0};
\node[mc, text=slate] at (6.12,-0.62) {0};
\node[mc, text=slate] at (6.80,-0.62) {0};
\node[lab, anchor=east] at (-0.5,-1.24) {after coin 3};
\node[mc, fill=accentsoft, draw=accent] at (0.00,-1.24) {1};
\node[mc, text=slate] at (0.68,-1.24) {0};
\node[mc, fill=accentsoft, draw=accent] at (1.36,-1.24) {1};
\node[mc, fill=amberbg, draw=amber, text=amber] at (2.04,-1.24) {1};
\node[mc, text=slate] at (2.72,-1.24) {0};
\node[mc, fill=amberbg, draw=amber, text=amber] at (3.40,-1.24) {1};
\node[mc, text=slate] at (4.08,-1.24) {0};
\node[mc, text=slate] at (4.76,-1.24) {0};
\node[mc, text=slate] at (5.44,-1.24) {0};
\node[mc, text=slate] at (6.12,-1.24) {0};
\node[mc, text=slate] at (6.80,-1.24) {0};
\node[lab, anchor=east] at (-0.5,-1.86) {after coin 5};
\node[mc, fill=accentsoft, draw=accent] at (0.00,-1.86) {1};
\node[mc, text=slate] at (0.68,-1.86) {0};
\node[mc, fill=accentsoft, draw=accent] at (1.36,-1.86) {1};
\node[mc, fill=accentsoft, draw=accent] at (2.04,-1.86) {1};
\node[mc, text=slate] at (2.72,-1.86) {0};
\node[mc, fill=accentsoft, draw=accent] at (3.40,-1.86) {1};
\node[mc, text=slate] at (4.08,-1.86) {0};
\node[mc, fill=amberbg, draw=amber, text=amber] at (4.76,-1.86) {1};
\node[mc, fill=amberbg, draw=amber, text=amber] at (5.44,-1.86) {1};
\node[mc, text=slate] at (6.12,-1.86) {0};
\node[mc, fill=amberbg, draw=amber, text=amber] at (6.80,-1.86) {1};
\node[labA, anchor=west, align=left] at (8.0,-0.6) {\texttt{dp |= dp << c}\\ copies every set bit $c$ places up};
\node[labM, anchor=west] at (0,-2.8) {amber = newly reachable; final set $\{2,3,5,7,8,10\}$ (6 sums, bit $0$ excluded)};
\end{tikzpicture}

\caption{The shift-and-or trick. Notice each coin copies every
set bit $w$ positions up --- parallelism for free.}
\label{fig:bitset}
\end{figure}
Further up the band: convex-hull trick, divide-and-conquer DP,
Knuth optimization, SOS DP, profile DP --- each has its own
structural precondition; treat them as named routes (see Notes).

\begin{exercise}[easy]\label{exr:mask}
Iterate the nonempty submasks of mask $m$ in $O(2^{|m|})$ ---
write the loop and argue each visited value is a submask.
\end{exercise}
\begin{exercise}[medium]\label{exr:tight}
Why can't the digit-DP memo be shared across tight=1 states?
What invariant does the flag preserve?
\end{exercise}
\begin{exercise}[medium]\label{exr:bylen}
Argue by-length is the \emph{only} correct fill order for
dp[l][r] = f(dp[l+1][r], dp[l][r-1]).
\end{exercise}
\begin{exercise}[medium]\label{exr:rerootcarries}
In the all-roots matching problem, what exactly does the second
pass carry from parent to child?
\end{exercise}

\section{Chapter summary}
\begin{table}[ht]\centering
{\small
\begin{tabularx}{\linewidth}{@{}lYYY@{}}
\toprule
\tablehead\thead{State} & \thead{Template} & \thead{Size} & \thead{Watch}\\
\midrule
subset & mask DP & $O(2^n n)$ & $n\le20$\\
(position, flags) & digit DP & $O(\log n)$ & tight not memoizable\\
interval & range DP & $O(n^3)$ worst & by-length only\\
(vertex, in/out) & tree DP & $O(n)$ & state pair = interface\\
sum-set & bitset & $O(nC/64)$ & word width\\
\bottomrule
\end{tabularx}}
\caption{Pick the state shape, then the order that respects its
dependency lattice.}
\label{tab:dp2}
\end{table}

\section{Common mistakes}
\index{common mistakes}
\begin{itemize}
\item \textbf{Memoizing tight.} The tight state depends on $N$'s
prefix --- caching it leaks between queries; memoize only loose.
\item \textbf{Wrong mask order.} Adding elements needs masks in any
order that adds bits; removing needs the reverse --- mixing
directions double-counts.
\item \textbf{Range filled by $l$.} dp[l][r] must see smaller
ranges first; any order not sorted by $r-l$ reads stale cells.
\item \textbf{Forgetting leading-zero.} ``Number not started'' is a
distinct state --- counting without it counts padded numbers.
\item \textbf{Bitmask index off-by-one.} Bit $i$ means element $i$,
not $i{+}1$ --- the bug shows only at $n$ small enough to enumerate.
\end{itemize}

\section{Additional exercises}
\begin{exercise}[easy] AC \emph{Matching} --- assignment mask DP. \end{exercise}
\begin{exercise}[easy] AC \emph{Grouping} --- partition mask DP. \end{exercise}
\begin{exercise}[easy] CSES \emph{Edit Distance}. \end{exercise}
\begin{exercise}[easy] USACO-G \emph{Barn Painting} --- tree counting DP. \end{exercise}
\begin{exercise}[medium] CSES \emph{Elevator Rides}. \end{exercise}
\begin{exercise}[medium] CSES \emph{Rectangle Cutting}. \end{exercise}
\begin{exercise}[medium] CSES \emph{Minimal Grid Path} --- lexicographic DP. \end{exercise}
\begin{exercise}[medium] USACO-G \emph{Uddered but not Herd} --- bitmask ordering. \end{exercise}
\begin{exercise}[medium] AC \emph{Independent Set} --- tree in/out. \end{exercise}
\begin{exercise}[medium] CF \emph{Team Building} --- mask + DP. \end{exercise}
\begin{exercise}[hard] CSES \emph{Counting Tilings} --- profile DP preview. \end{exercise}
\begin{exercise}[hard] CSES \emph{Counting Numbers}. \end{exercise}
\begin{exercise}[hard] CSES \emph{Longest Common Subsequence}. \end{exercise}
\begin{exercise}[hard] USACO-G \emph{Lights Off} --- mask + range DP. \end{exercise}

\chapternotes
The mask/digit/range templates follow Laaksonen's DP
chapters~\cite{laaksonen2024} and Sannemo~\cite{sannemo2025}; digit
DP's tight automaton and the bitset trick are cp-algorithms and the
USACO Guide's treatment~\cite{usacoguide,cpalgorithms}; the
optimization horizon is signposted in Halim~\cite{halimcp4}; archive
problems are official CSES/USACO~\cite{cses}.
